% Source : cours/cm2.qmd. Élasticité et incidence ; pas de donnée inventée.
\begin{board}{2}{Tabac : les achats réagissent au prix}
 {Ventes légales $\neq$ consommation : achats à l'étranger et autres circuits.}
\stage{5.4}{1}{Le paquet renchérit}
\stage{16}{2}{Une réaction qui évolue}
\stage{26.6}{3}{Versement et charge}
\packet{3.7}{8.1}{amber}
\packet{5}{8.1}{amber}
\bill{7.1}{8}
\draw[flow] (6.9,10.6)--(6.9,12);
\captext{5.4}{Une taxe augmente\\le prix payé.}
\noteat{16}{12.2}{Court terme}
\foreach \x in {12.9,14.1,15.3,16.5,17.7}{\packet{\x}{10}{blue}}
\noteat{16}{8.9}{Long terme}
\foreach \x in {14.1,15.3,16.5}{\packet{\x}{6.8}{blue}}
\captext{16}{L'effet du prix est\\plus fort à long terme.}
\shop{24.2}{8}
\bank{29}{8}
\coin{26.6}{10.9}
\draw[flow] (26,11.3)--(28,11.3);
\noteat{26.6}{12.3}{Le vendeur verse la taxe}
\node[smalllabel,text width=3cm] at (24.2,7.2) {Vendeur};
\node[smalllabel,text width=3cm] at (29,7.2) {État};
\captext{26.6}{Côté peu élastique :\\charge plus importante.}
\end{board}
